% =============================================================================
% Modèles de Markov Cachés
% =============================================================================

\minitoc

Ce chapitre traite des \emph{Modèles de Markov Cachés} (Hidden Markov Models, HMM). 
Historiquement développés pour le traitement automatique de la parole (\emph{Automatic Speech Recognition}, Rabiner, 1989) et le traitement du langage naturel (étiquetage morpho-syntaxique \emph{POS tagging}, modèles de langage, alignement bilingue), les HMM offrent un cadre probabiliste rigoureux pour les séquences discrètes. Bien que les modèles de fondation neuronaux (Transformers) dominent aujourd'hui ces applications, les algorithmes sous-jacents aux HMM (Forward-Backward, Viterbi, Baum-Welch/EM) constituent les bases fondamentales des modèles graphiques probabilistes et de la programmation dynamique.

% =============================================================================
\section{Modèles de Langage et Hypothèses Markoviennes}

Dans ce cours, on introduit les HMM en tant que modèles de langage. On comment ici par formaliser mathématiquement ces modèles en utilisant les chaînes de Markov. 

\subsection{Modélisation séquentielle par règle de la chaîne}

Soit $\mathcal{V}$ un vocabulaire fini de symboles (mots ou \emph{tokens}). On cherche à définir une loi de probabilité conjointe sur des séquences de longueur arbitraire bordées par deux balises déterministes $\operatorname{start}$ et $\operatorname{stop}$ :
\begin{equation}
    O = (\operatorname{start}, O_1, O_2, \dots, O_n, \operatorname{stop}) \in \mathcal{V}^*.
\end{equation}

Par application exacte de la règle de décomposition en chaîne des probabilités conditionnelles, toute distribution sur une séquence s'écrit de manière causale (contexte gauche) :
\begin{equation}
    \myP(\operatorname{start}, O_1, \dots, O_n, \operatorname{stop}) = \prod_{i=1}^{n+1} \myP(O_i \mid \operatorname{start}, O_1, \dots, O_{i-1})
\end{equation}
en posant conventionnellement $O_{n+1} = \operatorname{stop}$.

\begin{remark}[Intuition des dés de Noah Smith et explosion combinatoire]
    On peut interpréter ce processus génératif à l'aide d'une collection de dés dont les faces représentent les mots de $\mathcal{V}$. À chaque étape, l'historique complet $h = (\operatorname{start}, O_1, \dots, O_{i-1})$ détermine le dé à lancer, dont la distribution sur les faces fournit le mot suivant $O_i$. 
    Bien que d'une expressivité totale, ce modèle présente deux écueils rédhibitoires :
    \begin{enumerate}
        \item \textbf{Explosion paramétrique} : pour un historique de longueur au plus $n$ et un vocabulaire de taille $V = |\mathcal{V}|$, le nombre d'historiques possibles vaut $\sum_{k=0}^{n-1} V^k = \frac{V^n - 1}{V - 1}$. Pour $V = 10^5$ et $n = 10$, cela exigerait de stocker les paramètres d'environ $10^{50}$ distributions catégorielles, ce qui est matériellement impossible.
        \item \textbf{Défaut de généralisation} : toute séquence contenant un préfixe non observé lors de l'apprentissage se voit affecter une probabilité nulle.
    \end{enumerate}
\end{remark}

% -----------------------------------------------------
\subsection{Du Sac de Mots au Modèle de Markov d'ordre $m$}

Pour rendre l'estimation calculable, on réduit l'historique conditionnel via des hypothèses d'indépendance conditionnelle.

\begin{definition}[Modèle Sac de Mots / Unigramme]
    L'hypothèse la plus forte pose l'indépendance mutuelle totale de chaque observation vis-à-vis du contexte passé :
    \begin{equation}
        \myP(\operatorname{start}, O_1, \dots, O_n, \operatorname{stop}) = \prod_{i=1}^{n} \myP(O_i).
    \end{equation}
    Ce modèle ne requiert que $V - 1$ paramètres libres. En revanche, il détruit toute syntaxe : la probabilité d'une séquence est invariante par permutation de ses termes.
\end{definition}

Le compromis canonique consiste à borner la mémoire du processus aux $m$ derniers symboles émis.

\begin{definition}[Chaîne de Markov d'ordre $m$ / $n$-grammes]
    Un modèle de Markov stationnaire d'ordre $m \geqslant 1$ pose que la loi conditionnelle du token courant ne dépend que des $m$ tokens immédiatement antérieurs :
    \begin{equation}
        \myP(O_i \mid \operatorname{start}, O_1, \dots, O_{i-1}) = \myP(O_i \mid O_{i-m}, \dots, O_{i-1}).
    \end{equation}
    La probabilité conjointe de la séquence s'écrit alors :
    \begin{equation}
        \myP(\operatorname{start}, O_1, \dots, O_n, \operatorname{stop}) = \prod_{i=1}^{n+1} \myP(O_i \mid O_{i-m}^{i-1})
    \end{equation}
    où $O_{i-m}^{i-1} = (O_{i-m}, \dots, O_{i-1})$.
\end{definition}

\begin{remark}[Taxonomie et complexité des $n$-grammes]
    Dans la terminologie du traitement automatique du langage :
    \begin{itemize}
        \item $m = 0$ correspond au modèle \emph{unigramme} ($V-1$ paramètres) ;
        \item $m = 1$ correspond au modèle \emph{bigramme} ($V(V-1)$ paramètres de transition) ;
        \item $m = 2$ correspond au modèle \emph{trigramme} ($V^2(V-1)$ paramètres).
    \end{itemize}
    Au-delà de $m = 3$ ou $4$, la matrice des comptes devient extrêmement creuse (\emph{sparsity}), rendant indispensables des techniques de lissage statistique (Good-Turing, Kneser-Ney).
\end{remark}

% -----------------------------------------------------
\subsection{Modèles à classes de Brown}

Pour injecter de la régularité syntaxique sans subir l'explosion combinatoire des $n$-grammes lexicaux, Brown et al. (1992) partitionnent le vocabulaire $\mathcal{V}$ en $C \ll V$ classes syntaxiques disjointes via une application déterministe $\operatorname{cl} : \mathcal{V} \to \llbracket 1, C \rrbracket$ obtenue par partitionnement statistique :

\begin{definition}[Modèle de séquence fondé sur des classes]
    La génération d'une observation $O_i$ est scindée en deux lois : une transition d'état à état entre classes, suivie de l'émission du mot conditionnellement à sa propre classe :
    \begin{equation}
        \myP(\operatorname{start}, O_1, \dots, O_n, \operatorname{stop}) = \prod_{i=1}^{n+1} \myP(O_i \mid \operatorname{cl}(O_i)) \cdot \myP(\operatorname{cl}(O_i) \mid \operatorname{cl}(O_{i-1})).
    \end{equation}
\end{definition}

En notant $A \in \R^{C \times C}$ la matrice de transition inter-classes et $B \in \R^{C \times V}$ la matrice d'émission des mots sachant la classe, ce modèle ne requiert que $C^2 + C \times V$ paramètres. 

Cependant, ce modèle suppose que la fonction de projection $\operatorname{cl}$ est déterministe et connue à l'avance. Or, dans la langue naturelle, un mot est fondamentalement polysémique : le mot « \emph{flies} » peut être un verbe ou un nom selon le contexte. L'état syntaxique sous-jacent est donc par essence \emph{non observé} : c'est l'acte de naissance des \emph{Modèles de Markov Cachés}.

% =============================================================================
\section{Formalisme des Modèles de Markov Cachés (HMM)}

\subsection{Définition mathématique}

Un modèle de Markov caché de premier ordre modélise un processus bivarié $(S_t, O_t)_{t \ge 1}$ où la suite des états $(S_t)$ forme une chaîne de Markov non observable (latente), et $(O_t)$ est la suite des observations visibles émises à chaque instant conditionnellement à l'état courant.

\begin{definition}[Paramètres canoniques d'un HMM]
    Un HMM à états discrets est formellement caractérisé par le triplet $\lambda = (\pi, A, B)$ :
    \begin{enumerate}
        \item Un ensemble fini d'états cachés $\mathcal{S} = \{s_1, \dots, s_N\}$.
        \item Un alphabet fini d'observations $\mathcal{V} = \{v_1, \dots, v_M\}$.
        \item La distribution initiale des états $\pi = (\pi_i)_{i=1}^N$ avec $\pi_i = \myP(S_1 = s_i)$.
        \item La matrice de transition $A = (a_{ij}) \in \R^{N \times N}$ régie par l'hypothèse markovienne :
        \begin{equation}
            a_{ij} = \myP(S_{t+1} = s_j \mid S_t = s_i), \quad \sum_{j=1}^N a_{ij} = 1.
        \end{equation}
        \item La matrice d'émission $B = (b_j(k)) \in \R^{N \times M}$ vérifiant l'hypothèse d'indépendance conditionnelle :
        \begin{equation}
            b_j(k) = \myP(O_t = v_k \mid S_t = s_j), \quad \sum_{k=1}^M b_j(k) = 1.
        \end{equation}
    \end{enumerate}
\end{definition}

Pour une séquence d'états $S = (S_1, \dots, S_T)$ et d'observations $O = (O_1, \dots, O_T)$, la loi conjointe se factorise immédiatement grâce aux indépendances conditionnelles du graphe :
\begin{equation}
    \myP(O, S \mid \lambda) = \pi_{S_1} b_{S_1}(O_1) \prod_{t=2}^T a_{S_{t-1}, S_t} b_{S_t}(O_t).
\end{equation}

\subsection{Les trois problèmes fondamentaux de Rabiner (1989)}

L'utilisation des modèles HMM induit trois problèmes fondamentaux appelés des \emph{problèmes de Rabiner}, en l'honneur de la personne qui les a formalisés la première fois en 1989. 

\begin{proposition}[Problèmes de Rabiner - 1989]
    Les 3 problèmes présentés par Rabiner sont : 
    \begin{enumerate}
        \item \textbf{L'Évaluation} : À partir d'une séquence d'observation donnée $O = (O_1, \dots, O_m)$ comment calculer la probabilité qu'elle soit générée par le modèle ? 
        L'approche naïve serait de dénombrer toutes les séquences possibles de probabilité non nulle, mais c'est un problème de complexité exponentielle par rapport à la taille de la séquence. En pratique, on utilisera l'algorithme de programmation dynnamique \emph{Forward}. 
        \item \textbf{Inférence/Décodage} : À partir d'une observation $O = (O_1, \dots, O_n)$ quelle est la séquence d'états cachés $S = (S_1, \dots, S_n)$ qui aurait généré ces observations la plus probable ? 
        On pourrait aussi calculer des fréquences sur toutes les séquences d'états possibles pouvant générer cette observation mais c'est aussi de complexité exponentiell. On utilisera l'algorithme de \emph{Viterbi}. 
        \item \textbf{L'Apprentissage} : À partir d'un jeu de données d'observations, comment estimer les matrices de transition $A$ et $B$ pour construire le modèle ? On distingue ici deux cas : 
        \begin{itemize}
            \item \textbf{Apprentissage Supervisé} : On dispose ici des séquences d'observations et les états cachés correspondant : il suffit alors de calculer des fréquences dessus pour estimer les deux matrices. 
            \item \textbf{Apprentissage Non-Supervisé} : On ne dispose que des séquences d'observations. On utilisera pour cela l'algorithme \emph{Expectation-Maximization (EM)} développé par Dempster en 1977. 
            Cet algorithme se base sur le calcul d'un maximum de vraissemblance sur les états possibles. 

            On peut aussi citer des approches par clustering sur les observations pour définir des états cachés ou des $k$-means. 
        \end{itemize}
    \end{enumerate}
\end{proposition}

% =============================================================================
% Résolution Algorithmique des Trois Problèmes de Rabiner

\section{Résolution Algorithmique des Trois Problèmes de Rabiner}

% -----------------------------------------------------------------------------
\subsection{Problème 1 : Évaluation et Algorithme Forward-Backward}

Soit un modèle HMM fixé $\lambda = (\pi, A, B)$ et une séquence d'observations $O = (O_1, O_2, \dots, O_T) \in \mathcal{V}^T$. Le problème d'évaluation consiste à déterminer la vraisemblance marginale $\myP(O \mid \lambda)$. Commençons par étudier la complexité de l'approche naïve pour se convaincre de l'utilisation d'algorithmes plus performants. 

\begin{proposition}[Approche naïve]
    Par la loi des probabilités totales, la vraisemblance marginale s'obtient en sommant la loi conjointe sur l'ensemble $\mathcal{S}^T$ des trajectoires d'états latents possibles :
    \begin{equation}
        \myP(O \mid \lambda) = \sum_{S \in \mathcal{S}^T} \myP(O, S \mid \lambda) = \sum_{s_{i_1}=1}^N \dots \sum_{s_{i_T}=1}^N \pi_{s_{i_1}} b_{s_{i_1}}(O_1) \prod_{t=2}^T a_{s_{i_{t-1}}, s_{i_t}} b_{s_{i_t}}(O_t).
    \end{equation}
    L'espace $\mathcal{S}^T$ possède un cardinal de $N^T$. Chaque terme nécessitant $2T-1$ multiplications, la complexité temporelle brute est en $\mathcal{O}(T \cdot N^T)$. Ce calcul devient physiquement irréalisable dès que $T$ dépasse quelques dizaines d'instants (pour $N=10$ et $T=50$, $N^T = 10^{50}$).    
\end{proposition}

On préférera donc l'algorithme Forward ou Backward définis ci-dessous. L'algorithme Forward résout ce problème par \emph{programmation dynamique} en factorisant les calculs sur un treillis d'états (\emph{trellis}).

\begin{proposition}[Algorithme Forward]
    On introduit la variable intermédiaire $\alpha_t(i)$, définie comme la probabilité conjointe du préfixe d'observations jusqu'à l'instant $t$ et de l'état latent au temps $t$ :
    \begin{equation}
        \alpha_t(i) = \myP(O_1, \dots, O_t, S_t = s_i \mid \lambda), \quad \forall i \in \llbracket 1, N \rrbracket, \, t \in \llbracket 1, T \rrbracket.
    \end{equation}

    \begin{algorithm}[H]
    \caption{Algorithme Forward}
    \begin{enumerate}
        \item \textbf{Initialisation} ($t = 1$) :
        \begin{equation}
            \alpha_1(i) = \pi_i \, b_i(O_1), \quad \forall i \in \llbracket 1, N \rrbracket.
        \end{equation}
        \item \textbf{Induction / Récurrence} ($t = 1, \dots, T-1$) :
        \begin{equation}
            \alpha_{t+1}(j) = \left[ \sum_{i=1}^N \alpha_t(i) \, a_{ij} \right] b_j(O_{t+1}), \quad \forall j \in \llbracket 1, N \rrbracket.
        \end{equation}
        \item \textbf{Terminaison} :
        \begin{equation}
            \myP(O \mid \lambda) = \sum_{i=1}^N \alpha_T(i).
        \end{equation}
    \end{enumerate}
    \end{algorithm}
\end{proposition}

\begin{remark}[Complexité]
    Le treillis compte $N \times T$ nœuds. À chaque transition temporelle $t \to t+1$, le calcul de $\alpha_{t+1}(j)$ agrège les contributions des $N$ états antérieurs via $N$ multiplications et additions. La complexité temporelle chute ainsi à $\mathcal{O}(N^2 T)$, avec une complexité spatiale en mémoire de $\mathcal{O}(N T)$ (ou $\mathcal{O}(N)$ si l'on ne conserve que le vecteur de l'instant précédent).
\end{remark}

\begin{proposition}[Algorithme Backward]
    De manière duale, on définit la variable Backward $\beta_t(i)$ comme la probabilité conditionnelle du suffixe d'observations allant de $t+1$ à $T$, conditionnellement à l'état latent présent $S_t = s_i$ :
    \begin{equation}
        \beta_t(i) = \myP(O_{t+1}, \dots, O_T \mid S_t = s_i, \lambda), \quad \forall i \in \llbracket 1, N \rrbracket, \, t \in \llbracket 1, T \rrbracket.
    \end{equation}

        \begin{algorithm}[H]
        \caption{Algorithme Backward}
        \begin{enumerate}
            \item \textbf{Initialisation} ($t = T$) :
            \begin{equation}
                \beta_T(i) = 1, \quad \forall i \in \llbracket 1, N \rrbracket.
            \end{equation}
            \item \textbf{Induction régressive} ($t = T-1, \dots, 1$) :
            \begin{equation}
                \beta_t(i) = \sum_{j=1}^N a_{ij} \, b_j(O_{t+1}) \, \beta_{t+1}(j), \quad \forall i \in \llbracket 1, N \rrbracket.
            \end{equation}
            \item \textbf{Terminaison} :
            \begin{equation}
                \myP(O \mid \lambda) = \sum_{i=1}^N \pi_i \, b_i(O_1) \, \beta_1(i).
            \end{equation}
        \end{enumerate}
    \end{algorithm}
\end{proposition}

% -----------------------------------------------------------------------------
\subsection{Problème 2 : Décodage et Inférence}

Étant donnée une séquence observée $O$, le décodage consiste à inférer la séquence d'états latents sous-jacente. Selon l'objectif visé, deux formulations distinctes sont envisagées.
Le critère le plus classique recherche la trajectoire complète $S^* = (S_1^*, \dots, S_T^*)$ qui maximise la loi conjointe (et donc la loi a posteriori globale $\myP(S \mid O, \lambda)$) :
\begin{equation}
    S^* = \arg\max_{S \in \mathcal{S}^T} \myP(S \mid O, \lambda) = \arg\max_{S \in \mathcal{S}^T} \myP(O, S \mid \lambda).
\end{equation}
On introduit la variable de Viterbi $\delta_t(i)$, représentant la probabilité maximale d'une séquence partielle d'états se terminant en $s_i$ au temps $t$, ainsi qu'une matrice de pointeurs $\psi_t(i)$ (\emph{backpointers}) stockant l'antécédent optimal :
\begin{equation}
    \delta_t(i) \triangleq \max_{S_1, \dots, S_{t-1}} \myP(O_1, \dots, O_t, S_1, \dots, S_{t-1}, S_t = s_i \mid \lambda).
\end{equation}

\begin{algorithm}[H]
\caption{Algorithme de Viterbi}
\begin{enumerate}
    \item \textbf{Initialisation} ($t = 1$) :
    \begin{align}
        \delta_1(i) &= \pi_i \, b_i(O_1), \quad \forall i \in \llbracket 1, N \rrbracket, \quad 
        \psi_1(i) &= 0.
    \end{align}
    \item \textbf{Récurrence} ($t = 2, \dots, T$) :
    \begin{align}
        \delta_t(j) &= \left[ \max_{1 \le i \le N} \big( \delta_{t-1}(i) \, a_{ij} \big) \right] b_j(O_t), \quad \forall j \in \llbracket 1, N \rrbracket, \\
        \psi_t(j) &= \arg\max_{1 \le i \le N} \big( \delta_{t-1}(i) \, a_{ij} \big), \quad \forall j \in \llbracket 1, N \rrbracket.
    \end{align}
    \item \textbf{Terminaison} :
    \begin{align}
        P^* &= \max_{1 \le i \le N} \delta_T(i), \\
        S_T^* &= \arg\max_{1 \le i \le N} \delta_T(i).
    \end{align}
    \item \textbf{Reconstruction du chemin (\emph{Backtracking})} ($t = T-1, \dots, 1$) :
    \begin{equation}
        S_t^* = \psi_{t+1}(S_{t+1}^*).
    \end{equation}
\end{enumerate}
\end{algorithm}

\begin{proposition}[\emph{Posterior Decoding}]
    Une alternative consiste à déterminer, pour chaque pas de temps $t$ pris isolément, l'état qui maximise la probabilité a posteriori conditionnelle $\myP(S_t = s_i \mid O, \lambda)$ :
    \begin{equation}
        S_*^t = \arg\max_{1 \le i \le N} \myP(S_t = s_i \mid O, \lambda).
    \end{equation}
    Grâce aux indépendances conditionnelles du modèle, la probabilité conjointe de l'état $S_t = s_i$ et de l'ensemble de la séquence d'observations $O$ se factorise par le produit des variables Forward et Backward :
    \begin{equation}
        \myP(S_t = s_i, O \mid \lambda) = \alpha_t(i) \, \beta_t(i).
    \end{equation}
    Par la formule de Bayes :
    \begin{equation}
        \myP(S_t = s_i \mid O, \lambda) = \frac{\myP(S_t = s_i, O \mid \lambda)}{\myP(O \mid \lambda)} = \frac{\alpha_t(i) \, \beta_t(i)}{\sum_{j=1}^N \alpha_t(j) \, \beta_t(j)}.
    \end{equation}
    Le décodage a posteriori revient donc à évaluer :
    \begin{equation}
        S_*^t = \arg\max_{1 \le i \le N} \big( \alpha_t(i) \, \beta_t(i) \big).
    \end{equation}
\end{proposition}

\begin{remark}[Comparaison Viterbi vs Décodage a posteriori]
    \begin{itemize}
        \item \textbf{Viterbi} garantit la validité structurelle globale du chemin prédit : aucune transition de probabilité nulle ($a_{ij} = 0$) ne peut être empruntée.
        \item \textbf{Le décodage a posteriori} minimise le taux d'erreur moyen par position temporelle. Cependant, en maximisant chaque état indépendamment de ses voisins, il peut produire deux états consécutifs $S_*^t$ et $S_*^{t+1}$ pour lesquels la transition est impossible ($a_{S_*^t, S_*^{t+1}} = 0$). Dans ce cas, la probabilité globale de la séquence d'états résultante est strictement nulle.
    \end{itemize}
\end{remark}

Les algorithmes Forward et Viterbi partagent la même topologie de programmation dynamique sur le treillis. Ils correspondent à une instanciation formelle du même algorithme abstrait sur deux \emph{semi-anneaux} distincts $(K, \oplus, \otimes, \bar{0}, \bar{1})$ :

\begin{table}[H]
    \centering
    \begin{tabularx}{\textwidth}{l X X}
    \hline
    \textbf{Caractéristique} & \textbf{Semi-anneau des Réels (Forward)} & \textbf{Semi-anneau Tropical (Viterbi)} \\
    \hline
    Support $K$ & $\R_+$ & $\R_+$ \\
    Élément nul $\bar{0}$ & $0$ & $0$ \\
    Élément unité $\bar{1}$ & $1$ & $1$ \\
    Addition $\oplus$ & $+$ (somme des alternatives) & $\max$ (sélection du meilleur chemin) \\
    Multiplication $\otimes$ & $\times$ (accumulation causale) & $\times$ (accumulation causale) \\
    \hline
    \end{tabularx}
    \caption{Correspondance algébrique entre Forward et Viterbi via les semi-anneaux.}
\end{table}

% -----------------------------------------------------------------------------
\subsection{Stabilité Numérique et Domaine Logarithmique}

Sur de longues séquences, la multiplication itérative de probabilités inférieures à $1$ entraîne rapidement un dépassement de capacité par le bas (\emph{arithmetic underflow}) sur les flottants standard.

Pour éviter l'écrasement à zéro, on manipule les logarithmes des probabilités. Pour deux log-probabilités $u = \ln p$ et $v = \ln q$ :
\begin{itemize}
    \item La multiplication devient une addition : $\ln(p \times q) = u + v$.
    \item Dans l'algorithme de Viterbi, la monotonie du logarithme commute avec le maximum :
    \begin{equation}
        \ln \delta_t(j) = \ln b_j(O_t) + \max_{1 \le i \le N} \big[ \ln \delta_{t-1}(i) + \ln a_{ij} \big].
    \end{equation}
\end{itemize}

Pour l'algorithme Forward, l'addition de probabilités dans le domaine log nécessite le calcul de $\ln(p + q) = \ln(\mathrm{e}^u + \mathrm{e}^v)$. De manière générale, pour un vecteur $l = (l_1, \dots, l_n) \in \R^n$, on définit l'opérateur \emph{logsumexp} :
\begin{equation}
    \operatorname{LSE}(l) \triangleq \ln \left( \sum_{i=1}^n \mathrm{e}^{l_i} \right).
\end{equation}
Si certains composants de $l$ sont très grands ($l_i \gg 0$), le calcul brut de $\mathrm{e}^{l_i}$ provoque un dépassement par le haut (\emph{overflow}).

\begin{proposition}[Factorisation stable de Log-Sum-Exp]
    Pour tout scalaire $a \in \R$ :
    \begin{equation}
        \operatorname{LSE}(l) = a + \operatorname{LSE}(l - a) = a + \ln \left( \sum_{i=1}^n \mathrm{e}^{l_i - a} \right).
    \end{equation}
\end{proposition}

\begin{proof}
    En factorisant par $\mathrm{e}^a$ dans la somme :
    \begin{align}
        \operatorname{LSE}(l) &= \ln \left( \sum_{i=1}^n \mathrm{e}^{l_i} \right) = \ln \left( \sum_{i=1}^n \mathrm{e}^a \cdot \mathrm{e}^{l_i - a} \right) \nonumber \\
        &= \ln \left( \mathrm{e}^a \sum_{i=1}^n \mathrm{e}^{l_i - a} \right) = \ln(\mathrm{e}^a) + \ln \left( \sum_{i=1}^n \mathrm{e}^{l_i - a} \right) \nonumber \\
        &= a + \ln \left( \sum_{i=1}^n \mathrm{e}^{l_i - a} \right).
    \end{align}
\end{proof}

En pratique, on pose $a = \max_{1 \le i \le n} l_i$. Ainsi, pour tout $i$, $l_i - a \le 0$, ce qui borne $\mathrm{e}^{l_i - a} \in ]0, 1]$ et neutralise tout risque d'overflow.

% -----------------------------------------------------------------------------
\subsection{Problème 3 : Apprentissage des Paramètres}

L'objectif de l'apprentissage est d'estimer les paramètres du modèle $\lambda = (\pi, A, B)$ à partir d'un ensemble de séquences d'entraînement.

\begin{proposition}[Cas Supervisé]
    Lorsque les séquences d'observations et les séquences d'états associées sont intégralement observées (par exemple en étiquetage POS annoté manuellement), l'estimation des paramètres par \emph{Maximum de Vraisemblance} (MLE) se réduit à des ratios de fréquences empiriques :
    \begin{align}
        \hat{\pi}_i &= \frac{\operatorname{Count}(S_1 = s_i)}{\sum_{j=1}^N \operatorname{Count}(S_1 = s_j)}, \\
        \hat{a}_{ij} &= \frac{\operatorname{Count}(S_{t-1} = s_i, S_t = s_j)}{\sum_{k=1}^N \operatorname{Count}(S_{t-1} = s_i, S_t = s_k)}, \\
        \hat{b}_j(k) &= \frac{\operatorname{Count}(S_t = s_j, O_t = v_k)}{\sum_{m=1}^M \operatorname{Count}(S_t = s_j, O_t = v_m)}.
    \end{align}
\end{proposition}

\begin{proposition}[Cas Non-Supervisé : Algorithme EM]
    Lorsque les états latents ne sont pas observés, l'optimisation directe de la log-vraisemblance marginale $\ln \myP(O \mid \lambda) = \ln \sum_S \myP(O, S \mid \lambda)$ est non convexe. On applique l'algorithme d'\emph{Espérance-Maximisation} (EM), qui prend le nom d'algorithme de \emph{Baum-Welch} dans le cadre des HMM.

    L'algorithme itère deux étapes jusqu'à convergence :
    \begin{enumerate}
        \item \textbf{Étape E (Espérance)} : À l'aide des paramètres courants $\lambda^{(k)}$ et via les algorithmes Forward et Backward, on calcule les distributions a posteriori sur les variables latentes :
        \begin{itemize}
            \item Probabilité d'être dans l'état $s_i$ au temps $t$ :
            \begin{equation}
                \gamma_t(i) = \myP(S_t = s_i \mid O, \lambda^{(k)}) = \frac{\alpha_t(i) \, \beta_t(i)}{\sum_{j=1}^N \alpha_t(j) \, \beta_t(j)}.
            \end{equation}
            \item Probabilité de transiter de $s_i$ à $s_j$ entre les instants $t$ et $t+1$ :
            \begin{equation}
                \xi_t(i, j) = \myP(S_t = s_i, S_{t+1} = s_j \mid O, \lambda^{(k)}) = \frac{\alpha_t(i) \, a_{ij} \, b_j(O_{t+1}) \, \beta_{t+1}(j)}{\myP(O \mid \lambda^{(k)})}.
            \end{equation}
        \end{itemize}

        \item \textbf{Étape M (Maximisation)} : On met à jour les paramètres $\lambda^{(k+1)}$ en maximisant l'espérance de la log-vraisemblance des données complètes :
        \begin{align}
            \pi_i^{(k+1)} &= \gamma_1(i), \\
            a_{ij}^{(k+1)} &= \frac{\sum_{t=1}^{T-1} \xi_t(i, j)}{\sum_{t=1}^{T-1} \gamma_t(i)}, \\
            b_j(m)^{(k+1)} &= \frac{\sum_{t=1}^T \gamma_t(j) \cdot \mathbb{I}(O_t = v_m)}{\sum_{t=1}^T \gamma_t(j)}.
        \end{align}
    \end{enumerate}
    L'algorithme garantit la croissance monotone de la vraisemblance marginale $\myP(O \mid \lambda)$ à chaque itération et converge vers un extremum local du paysage d'optimisation.
\end{proposition}